    % ==== Reinforcement Learning =============================================================================
    \chapter{Reinforcement Learning}

    % -------------------------------------------------------
    \section{Introduction to Reinforcement Learning}

    \begin{tcbitemize}[ skin=sectionraster ]

        \tcbitem[title=What is Reinforcement Learning?, raster multicolumn=4]
        Reinforcement Learning (RL) is a learning paradigm in which an
        \textbf{agent} interacts with an \textbf{environment} to maximise
        cumulative \textbf{reward}. Unlike supervised learning (labelled
        examples) or unsupervised learning (structure discovery), the agent
        receives only a scalar reward signal --- often delayed --- and must learn
        by \emph{trial and error}.\textsuperscript{\cite[][p.~323]{hurbansGrokkingArtificialIntelligence2020}}

        The core challenge is the \textbf{exploration--exploitation trade-off}:
        should the agent exploit its best-known strategy or explore to
        discover potentially better ones?

        \tcbitem[title=The Reward Hypothesis, raster multicolumn=2]
        \textbf{Reward hypothesis}: all goals can be expressed as the
        maximisation of a scalar cumulative reward signal.

        \begin{itemize}
            \item Reward encodes the \emph{objective}, not the \emph{strategy}
            \item Reward may be immediate or heavily delayed
            \item The agent is not told \emph{how} to act --- only \emph{how well}
        \end{itemize}
        \textsuperscript{\cite[][p.~324]{hurbansGrokkingArtificialIntelligence2020}}

        \tcbitem[title={Agent--Environment Interaction Loop}, raster multicolumn=6]
        At each discrete timestep $t$ the agent receives state $s_t$, selects
        action $a_t$, and the environment returns reward $r_{t+1}$ and next
        state $s_{t+1}$. This produces a \textbf{trajectory}
        $\tau = (s_0, a_0, r_1, s_1, a_1, r_2, \ldots)$.%
        \textsuperscript{\cite[][p.~325]{hurbansGrokkingArtificialIntelligence2020}}
        \tcblower
        \begin{center}
        \begin{tikzpicture}[
            box/.style={draw, rounded corners=4pt, minimum width=2.2cm,
                        minimum height=0.9cm, align=center, font=\small},
            arr/.style={-Stealth, thick}
        ]
            \node[box, fill=LimeGreen!25] (agent) {Agent};
            \node[box, fill=orange!20, right=4.5cm of agent] (env) {Environment};
            \draw[arr] (agent.north) to[bend left=30]
                node[above, font=\footnotesize]{action $a_t$} (env.north);
            \draw[arr] (env.south) to[bend left=30]
                node[below, font=\footnotesize]{state $s_{t+1}$, reward $r_{t+1}$}
                (agent.south);
        \end{tikzpicture}
\captionof{figure}{Introduction to Reinforcement Learning: Agent--Environment Interaction Loop.}
\label{fig:ch07-reinforcement-learning-inline-01}
        \end{center}

        \tcbitem[title=Core Components, raster multicolumn=3]
        \begin{description}[leftmargin=0.5cm, labelindent=0pt]
            \item[\textbf{Agent}] The learner and decision-maker
            \item[\textbf{Environment}] Everything external to the agent
            \item[\textbf{State $s$}] Representation of the current situation
            \item[\textbf{Action $a$}] A choice available to the agent
            \item[\textbf{Reward $r$}] Scalar feedback from the environment
            \item[\textbf{Policy $\pi$}] Mapping: state $\to$ action (or distribution)
            \item[\textbf{Value $V(s)$}] Expected cumulative future reward from $s$
            \item[\textbf{Model}] Internal model of environment dynamics (optional)
        \end{description}
        \textsuperscript{\cite[][p.~326]{hurbansGrokkingArtificialIntelligence2020}}

        \tcbitem[title={RL vs.\ Supervised vs.\ Unsupervised}, raster multicolumn=3]
        \begin{tblr}{
            colspec={X[l]X[l]X[l]},
            hlines, vlines,
            row{1}={font=\bfseries\small, bg=LimeGreen!20},
            row{2-Z}={font=\small},
        }
            Supervised & Unsupervised & Reinforcement \\
            Labelled data & Unlabelled data & Interactions \\
            Immediate answer & No feedback & Delayed reward \\
            Generalise & Discover structure & Maximise return \\
            No exploration & No exploration & Exploration essential \\
        \end{tblr}
\captionof{table}{Introduction to Reinforcement Learning: RL vs.\ Supervised vs.\ Unsupervised.}
\label{tab:ch07-reinforcement-learning-inline-01}
        \textsuperscript{\cite[][p.~323]{hurbansGrokkingArtificialIntelligence2020}}

    \end{tcbitemize}


    % -------------------------------------------------------
    \section{Markov Decision Processes}

    \begin{tcbitemize}[ skin=sectionraster ]

        \tcbitem[title=The Markov Property, raster multicolumn=3]
        A state $s_t$ satisfies the \textbf{Markov property} if:
        \begin{equation}
            P(s_{t+1} \mid s_t, a_t) = P(s_{t+1} \mid s_0, a_0, \ldots, s_t, a_t)
        \end{equation}
        The future depends only on the \emph{current} state and action, not
        on the full history --- the state is a \emph{sufficient statistic} of
        the past. This memoryless property enables tractable dynamic
        programming solutions.\textsuperscript{\cite[][p.~140]{barberBayesianReasoningMachine2012}}

        \tcbitem[title={MDP: Formal Definition}, raster multicolumn=3]
        A \textbf{Markov Decision Process} is the tuple
        $\mathcal{M} = (\mathcal{S}, \mathcal{A}, P, R, \gamma)$:
        \begin{itemize}
            \item $\mathcal{S}$: finite state space
            \item $\mathcal{A}$: finite action space
            \item $P(s' \mid s, a)$: transition probability
            \item $R(s, a)$: expected immediate reward
            \item $\gamma \in [0,1]$: discount factor
        \end{itemize}
        The agent's goal is to find policy $\pi^*$ that maximises expected
        discounted return from every starting
        state.\textsuperscript{\cite[][p.~141]{barberBayesianReasoningMachine2012}}

        \tcbitem[title={Cumulative Discounted Return $G_t$}, raster multicolumn=6]
        The \textbf{return} at time $t$ is the discounted sum of all future
        rewards:\textsuperscript{\cite[][p.~142]{barberBayesianReasoningMachine2012}}
        \tcblower
        \begin{equation}
            G_t = \sum_{k=0}^{\infty} \gamma^k\, r_{t+k+1}
                = r_{t+1} + \gamma r_{t+2} + \gamma^2 r_{t+3} + \cdots
        \end{equation}
        Discount factor $\gamma$ controls the time-horizon:
        $\gamma = 0$ (myopic, only immediate reward);
        $\gamma \to 1$ (far-sighted, all future equally weighted).
        For $\gamma < 1$, the sum converges even for infinite horizons.

        \tcbitem[title=Value Functions, raster multicolumn=3]
        \textbf{State-value function} under policy $\pi$:
        \begin{equation}
            V^{\pi}(s) = \mathbb{E}_{\pi}\!\left[G_t \mid S_t = s\right]
        \end{equation}
        \textbf{Action-value (Q) function} under policy $\pi$:
        \begin{equation}
            Q^{\pi}(s, a) = \mathbb{E}_{\pi}\!\left[G_t \mid S_t = s, A_t = a\right]
        \end{equation}
        $V^{\pi}(s)$ is the expected return following $\pi$ from $s$;
        $Q^{\pi}(s,a)$ additionally fixes the first action.
        Optimal: $V^*(s) = \max_\pi V^\pi(s)$.\textsuperscript{\cite[][p.~143]{barberBayesianReasoningMachine2012}}

        \tcbitem[title=Bellman Equations, raster multicolumn=3]
        \textbf{Bellman expectation} decomposes the value function
        recursively:\textsuperscript{\cite[][p.~145]{barberBayesianReasoningMachine2012}}
        \begin{equation}
            V^{\pi}(s) = \sum_a \pi(a|s) \sum_{s'} P(s'|s,a)
                         \bigl[R(s,a) + \gamma V^{\pi}(s')\bigr]
        \end{equation}
        \textbf{Bellman optimality} characterises the optimal value:
        \begin{equation}
            V^*(s) = \max_a \sum_{s'} P(s'|s,a)
                     \bigl[R(s,a) + \gamma V^*(s')\bigr]
        \end{equation}

        \tcbitem[title=Value Iteration, raster multicolumn=3]
        A \textbf{dynamic programming} algorithm for finding $V^*$ when the
        model $(P, R)$ is
        known:\textsuperscript{\cite[][p.~153]{barberBayesianReasoningMachine2012}}
        \begin{enumerate}
            \item Initialise $V(s) \leftarrow 0\ \forall s$
            \item Repeat until $\|V_{k+1} - V_k\|_\infty < \epsilon$:
            \begin{equation}
                V_{k+1}(s) \leftarrow \max_a \sum_{s'} P(s'|s,a)
                           \bigl[R(s,a) + \gamma V_k(s')\bigr]
            \end{equation}
            \item Extract $\pi^*(s) = \arg\max_a \sum_{s'} P(s'|s,a)[R + \gamma V^*]$
        \end{enumerate}
        Convergence guaranteed (Bellman operator is a contraction).

        \tcbitem[title={MDP State Transitions}, raster multicolumn=3]
        MDP transitions depend on both state and action. Below: three states
        with two actions ($a_1, a_2$), illustrating stochastic transitions.
        \tcblower
        \begin{center}
        \begin{tikzpicture}[
            state/.style={circle, draw, minimum size=0.9cm, font=\small},
            arr/.style={-Stealth, thick},
            every node/.style={font=\small}
        ]
            \node[state, fill=LimeGreen!20] (s1) {$s_1$};
            \node[state, fill=orange!20, right=2.5cm of s1] (s2) {$s_2$};
            \node[state, fill=blue!15, below=1.5cm of s2] (s3) {$s_3$};
            \draw[arr, bend left=15] (s1) to
                node[above, font=\footnotesize]{$a_1,\,p{=}0.8$} (s2);
            \draw[arr] (s1) to
                node[left, font=\footnotesize]{$a_1,\,p{=}0.2$} (s3);
            \draw[arr] (s1) to[loop left]
                node[left, font=\footnotesize]{$a_2$} (s1);
            \draw[arr] (s2) to
                node[right, font=\footnotesize]{$a_1$} (s3);
            \draw[arr, bend right=20] (s3) to
                node[below, font=\footnotesize]{$a_2$} (s1);
        \end{tikzpicture}
\captionof{figure}{Markov Decision Processes: MDP State Transitions.}
\label{fig:ch07-reinforcement-learning-inline-02}
        \end{center}

        \tcbitem[title=Policy Types, raster multicolumn=3]
        \begin{description}[leftmargin=0.5cm]
            \item[\textbf{Deterministic}] $\pi : \mathcal{S} \to \mathcal{A}$
                --- one action per state
            \item[\textbf{Stochastic}] $\pi(a \mid s) = P(A{=}a \mid S{=}s)$
                --- probability distribution over actions
            \item[\textbf{Optimal $\pi^*$}] maximises $V^\pi(s)$ for every $s$;
                always exists as deterministic in finite MDPs
            \item[\textbf{Greedy}] $\pi(s) = \arg\max_a Q(s,a)$
                --- exploits current value estimate
            \item[\textbf{$\epsilon$-greedy}] greedy with probability $1{-}\epsilon$,
                random with probability $\epsilon$
        \end{description}
        \textsuperscript{\cite[][p.~144]{barberBayesianReasoningMachine2012}}

        \tcbitem[title=Markov Chain Monte Carlo (MCMC), raster multicolumn=3]
        \textbf{MCMC} uses Markov chains to draw samples from intractable
        distributions $p(x)$.

        \textbf{Metropolis-Hastings}:
        \begin{enumerate}
            \item Propose $x' \sim q(x' \mid x)$
            \item Compute acceptance ratio:
                  $\alpha = \min\!\left(1,\,
                  \frac{p(x')\,q(x \mid x')}{p(x)\,q(x' \mid x)}\right)$
            \item Accept $x'$ with probability $\alpha$; else stay at $x$
        \end{enumerate}
        At stationarity, samples follow $p(x)$.
        Applications: Bayesian posterior inference, policy
        evaluation.\textsuperscript{\cite[][p.~160]{barberBayesianReasoningMachine2012}}

    \end{tcbitemize}


    % -------------------------------------------------------
    \section{Bandit Problems}

    \begin{tcbitemize}[ skin=sectionraster ]

        \tcbitem[title={The Multi-Armed Bandit}, raster multicolumn=3]
        The \textbf{multi-armed bandit} (MAB) is an RL setting with
        \emph{no state transitions}: $k$ actions (arms), each yielding reward
        from an unknown distribution with mean $\mu_a$.

        Goal: minimise cumulative \textbf{regret}:
        \begin{equation}
            R_T = T\,\mu^* - \sum_{t=1}^{T} r_t
        \end{equation}
        where $\mu^* = \max_a \mu_a$.
        The fundamental tension is \textbf{exploration vs.\ exploitation}:
        exploit the best-known arm or explore to reduce
        uncertainty.\textsuperscript{\cite[][p.~335]{hurbansGrokkingArtificialIntelligence2020}}

        \tcbitem[title={$\epsilon$-Greedy Strategy}, raster multicolumn=3]
        With probability $\epsilon$ choose a random arm (explore); with
        probability $1 - \epsilon$ choose $\arg\max_a Q(a)$ (exploit).

        Incremental mean update after observing reward $r$:
        \begin{equation}
            Q(a) \leftarrow Q(a) + \frac{1}{n}\bigl[r - Q(a)\bigr]
        \end{equation}
        \begin{itemize}
            \item Simple and robust; $\epsilon$ often decayed over time
            \item Wastes exploration uniformly across all arms, even
                  well-understood ones
        \end{itemize}
        \textsuperscript{\cite[][p.~336]{hurbansGrokkingArtificialIntelligence2020}}

        \tcbitem[title={Upper Confidence Bound (UCB)}, raster multicolumn=3]
        Explore arms with high \emph{uncertainty} rather than randomly.
        \textbf{UCB1} selects:
        \begin{equation}
            A_t = \arg\max_a \left[Q_t(a) + c\,\sqrt{\frac{\ln t}{N_t(a)}}\right]
        \end{equation}
        where $N_t(a)$ is the number of times arm $a$ has been pulled.
        The confidence bonus decreases as an arm is explored more; $c$
        controls exploration width.
        UCB1 achieves $O(\ln T)$ regret --- asymptotically
        optimal.\textsuperscript{\cite[][p.~337]{hurbansGrokkingArtificialIntelligence2020}}

        \tcbitem[title=Thompson Sampling, raster multicolumn=3]
        A \textbf{Bayesian} approach: maintain a prior over arm reward
        parameters; sample from the posterior; pull the arm with the highest
        sampled value; update the posterior on the observed reward.

        For Bernoulli rewards, use a $\text{Beta}(\alpha_a, \beta_a)$ prior:
        \begin{itemize}
            \item Success: $\alpha_a \leftarrow \alpha_a + 1$
            \item Failure: $\beta_a \leftarrow \beta_a + 1$
        \end{itemize}
        Naturally balances exploration and exploitation via uncertainty.
        Empirically often outperforms $\epsilon$-greedy and
        UCB.\textsuperscript{\cite[][p.~338]{hurbansGrokkingArtificialIntelligence2020}}

    \end{tcbitemize}


    % -------------------------------------------------------
    \section{Q-Learning and Deep Reinforcement Learning}

    \begin{tcbitemize}[ skin=sectionraster ]

        \tcbitem[title=Temporal Difference Learning, raster multicolumn=3]
        \textbf{TD learning} is \emph{model-free}: learns directly from
        sampled transitions without a model of $P$ or $R$, and updates
        \emph{online} after each step.

        \textbf{TD(0)} value update:
        \begin{equation}
            V(s_t) \leftarrow V(s_t) + \alpha
            \underbrace{\bigl[r_{t+1} + \gamma V(s_{t+1}) - V(s_t)\bigr]}_{\text{TD error }\delta_t}
        \end{equation}
        TD(0) bootstraps from the current estimate; $\alpha$ is the learning
        rate.\textsuperscript{\cite[][p.~343]{hurbansGrokkingArtificialIntelligence2020}}

        \tcbitem[title={SARSA vs.\ Q-Learning}, raster multicolumn=3]
        Both learn $Q(s,a)$ via TD updates but differ in the \emph{target}:

        \textbf{SARSA} (on-policy):
        \begin{equation}
            Q(s,a) \leftarrow Q(s,a) + \alpha\bigl[r + \gamma Q(s',a') - Q(s,a)\bigr]
        \end{equation}
        Uses the \emph{actually taken} next action $a'$.

        \textbf{Q-Learning} (off-policy):
        \begin{equation}
            Q(s,a) \leftarrow Q(s,a) + \alpha\bigl[r + \gamma \max_{a'} Q(s',a') - Q(s,a)\bigr]
        \end{equation}
        Always bootstraps from the \emph{greedy} action; converges to
        $Q^*$ regardless of behaviour
        policy.\textsuperscript{\cite[][p.~346]{hurbansGrokkingArtificialIntelligence2020}}

        \tcbitem[title={Deep Q-Networks (DQN)}, raster multicolumn=6]
        For large or continuous state spaces, a Q-table is intractable.
        \textbf{DQN} approximates the Q-function with a deep neural network
        $Q(s,a;\theta) \approx Q^*(s,a)$, trained to minimise:
        \begin{equation}
            \mathcal{L}(\theta) = \mathbb{E}_{(s,a,r,s') \sim \mathcal{D}}
            \Bigl[\bigl(r + \gamma \max_{a'} Q(s',a';\theta^-) - Q(s,a;\theta)\bigr)^2\Bigr]
        \end{equation}
        Key innovations: (1)~\textbf{experience replay} buffer $\mathcal{D}$;
        (2)~\textbf{target network} $\theta^-$ updated periodically (not every
        step). Achieved human-level performance on 49 Atari
        games.\textsuperscript{\cite[][p.~350]{hurbansGrokkingArtificialIntelligence2020}}

        \tcbitem[title=Experience Replay, raster multicolumn=3]
        Store each transition $(s, a, r, s')$ in a circular \textbf{replay
        buffer} $\mathcal{D}$; sample random \emph{mini-batches} to train.

        Benefits:
        \begin{enumerate}
            \item \textbf{Decorrelates} temporally adjacent samples, reducing
                  variance
            \item \textbf{Sample efficiency}: each transition reused many times
            \item \textbf{Stabilises} gradient updates by breaking
                  non-stationarity
        \end{enumerate}
        \textsuperscript{\cite[][p.~351]{hurbansGrokkingArtificialIntelligence2020}}

        \tcbitem[title=Double Q-Learning, raster multicolumn=3]
        Standard Q-learning suffers from \textbf{overestimation bias}: the
        same network selects \emph{and} evaluates the maximum action, biasing
        targets upward.

        \textbf{Double DQN} decouples selection and evaluation:
        \begin{equation}
            y = r + \gamma\, Q\!\left(s',\, \arg\max_{a'} Q(s',a';\theta);\;\theta^-\right)
        \end{equation}
        Online network $\theta$ selects; target network $\theta^-$ evaluates.
        Reduces overestimation, yields more stable training and better
        policies.\textsuperscript{\cite[][p.~353]{hurbansGrokkingArtificialIntelligence2020}}

        \tcbitem[title={Sparse Rewards and Reward Shaping}, raster multicolumn=3]
        Many real tasks only give reward at task completion --- the agent faces
        a hard \textbf{credit assignment} problem.

        \textbf{Mitigation strategies}:
        \begin{itemize}
            \item \textbf{Reward shaping}: add potential-based intermediary
                  rewards $F(s,s') = \gamma \Phi(s') - \Phi(s)$ without
                  changing optimal policy
            \item \textbf{Intrinsic motivation}: curiosity bonus for novel states
            \item \textbf{Curriculum learning}: start with easier sub-tasks
            \item \textbf{Hindsight Experience Replay}: relabel failed
                  trajectories toward a different goal
        \end{itemize}
        \textsuperscript{\cite[][p.~357]{hurbansGrokkingArtificialIntelligence2020}}

        \tcbitem[title=Hierarchical Reinforcement Learning, raster multicolumn=3]
        Decomposes a long-horizon task into \textbf{sub-goals} and
        \textbf{sub-policies} at multiple levels of temporal abstraction.

        \textbf{Options framework}: an option is a triple
        $(I, \pi_\omega, \beta)$:
        \begin{itemize}
            \item $I \subseteq \mathcal{S}$: initiation set
            \item $\pi_\omega(a|s)$: intra-option policy
            \item $\beta(s)$: termination probability
        \end{itemize}
        A high-level policy selects options; each option executes its
        sub-policy until termination. Enables reuse of learned
        skills.\textsuperscript{\cite[][p.~358]{hurbansGrokkingArtificialIntelligence2020}}

        \tcbitem[title={Value-Based vs.\ Policy-Based Methods}, raster multicolumn=3]
        \begin{tblr}{
            colspec={X[l]X[l]},
            hlines, vlines,
            row{1}={font=\bfseries\small, bg=LimeGreen!20},
            row{2-Z}={font=\small},
        }
            Value-Based & Policy-Based \\
            Learn $Q$ or $V$; derive $\pi$ & Learn $\pi$ directly \\
            Q-learning, DQN & REINFORCE, PPO \\
            Natural for discrete actions & Handles continuous actions \\
            Deterministic output & Stochastic policy \\
            Can be unstable & High variance gradients \\
        \end{tblr}
\captionof{table}{Q-Learning and Deep Reinforcement Learning: Value-Based vs.\ Policy-Based Methods.}
\label{tab:ch07-reinforcement-learning-inline-02}
        \textsuperscript{\cite[][p.~360]{hurbansGrokkingArtificialIntelligence2020}}

        \tcbitem[title={Actor-Critic Methods}, raster multicolumn=3]
        Combines value-based and policy-based approaches:
        \begin{itemize}
            \item \textbf{Actor} $\pi(a|s;\theta)$: parametric policy, updated
                  via gradient ascent
            \item \textbf{Critic} $V(s; w)$: estimates state values, reduces
                  gradient variance
        \end{itemize}
        The \textbf{advantage function} $A(s,a) = Q(s,a) - V(s)$ measures
        how much better action $a$ is than average, reducing variance.

        \textbf{A2C}: synchronous workers; \textbf{A3C}: asynchronous
        parallel workers eliminating need for a replay
        buffer.\textsuperscript{\cite[][p.~362]{hurbansGrokkingArtificialIntelligence2020}}

    \end{tcbitemize}
    % -------------------------------------------------------
    \section{Reinforcement Learning Approaches}

    \subsection{Model-Free Learning}
    \begin{tcbitemize}[ skin=sectionraster ]

    \tcbitem[title=Model-Free RL, raster multicolumn=3]
    In \textbf{model-free} reinforcement learning the agent learns entirely through
    interaction with the environment --- no internal model of state transitions or
    rewards is built. The agent discovers which actions are valuable solely from
    trial-and-error experience.\textsuperscript{\cite[][p.~348]{hurbansGrokkingArtificialIntelligence2020}}

    Key characteristics:
    \begin{itemize}
        \item No explicit environment model required
        \item Learns from $(s, a, r, s’)$ tuples directly
        \item More data-hungry than model-based, but robust to model errors
        \item Examples: Q-learning, SARSA, REINFORCE
    \end{itemize}

    \tcbitem[title=The Q-Learning Update, raster multicolumn=3]
    \textbf{Q-learning} learns the action-value function $Q(s,a)$ --- the expected
    cumulative reward from taking action $a$ in state $s$ and following the
    optimal policy thereafter.\textsuperscript{\cite[][pp.~335--340]{hurbansGrokkingArtificialIntelligence2020}}

    The \textbf{Bellman update rule}:
    \begin{equation}
        Q(s,a) \leftarrow (1{-}\alpha)\,Q(s,a)
                 + \alpha\!\left[r + \gamma \max_{a’} Q(s’,a’)\right]
    \end{equation}
    $\alpha$: learning rate; $\gamma$: discount factor; $r$: immediate reward;
    $s’$: next state.

    The Q-table (rows = states, columns = actions) is initialised to zero and
    updated incrementally. At inference time: $a^* = \arg\max_a Q(s,a)$.

    \tcbitem[title={SARSA: On-Policy TD Control}, raster multicolumn=3]
    \textbf{SARSA} is \emph{on-policy}: it updates $Q$ using the action $a’$
    \emph{actually taken} by the current policy, not the greedy
    maximum.\textsuperscript{\cite[][Ch.~6]{suttonReinforcementLearningIntroduction2018}}

    \begin{equation}
        Q(s,a) \leftarrow Q(s,a) + \alpha\bigl[r + \gamma\,Q(s’,a’) - Q(s,a)\bigr]
    \end{equation}

    \begin{itemize}
        \item On-policy: more conservative, avoids dangerous exploratory actions
        \item Q-learning: faster convergence, may take risky shortcuts in training
        \item Both converge to the optimal $Q^*$ given sufficient exploration
    \end{itemize}

    \tcbitem[title=REINFORCE (Policy Gradient), raster multicolumn=3]
    \textbf{REINFORCE} directly optimises the policy $\pi_\theta(a|s)$ by
    gradient ascent on expected
    return.\textsuperscript{\cite[][Ch.~13]{suttonReinforcementLearningIntroduction2018}}

    \begin{equation}
        \nabla_\theta J(\theta) = \mathbb{E}_\pi\!\left[
            \nabla_\theta \ln \pi_\theta(a_t|s_t) \cdot G_t\right]
    \end{equation}

    where $G_t = \sum_{k=0}^{\infty} \gamma^k r_{t+k+1}$ is the return from step $t$.

    \begin{itemize}
        \item Monte Carlo: updates after complete episodes
        \item High variance; baseline subtraction (actor-critic) reduces variance
        \item Natural for continuous action spaces
    \end{itemize}

    \end{tcbitemize}

    \subsection{Model-Based Learning}
    \begin{tcbitemize}[ skin=sectionraster ]

    \tcbitem[title=Model-Based RL, raster multicolumn=3]
    In \textbf{model-based} RL the agent learns or is given an explicit model
    $\hat{M}(s,a) \to (s’, r)$ of the environment. It can \emph{plan} by
    simulating experience inside the model without executing every
    action.\textsuperscript{\cite[][p.~348]{hurbansGrokkingArtificialIntelligence2020}}

    \begin{itemize}
        \item \textbf{Sample efficient}: plan many steps from few real interactions
        \item Can reason about counterfactual actions before committing
        \item Model errors compound over long horizons --- accuracy matters
        \item Trade-off: model capacity vs.\ model bias
    \end{itemize}

    \tcbitem[title={Dyna-Q: Planning + Learning}, raster multicolumn=3]
    \textbf{Dyna-Q} (Sutton, 1991) combines real experience with simulated
    experience from a learned
    model:\textsuperscript{\cite[][Ch.~8]{suttonReinforcementLearningIntroduction2018}}

    \begin{enumerate}
        \item Take action $a$; observe $r$, $s’$ from real environment
        \item Update Q-table with real transition
        \item Update model: $\hat{M}(s,a) \leftarrow (s’,r)$
        \item Repeat $n$ times: sample random past $(s,a)$, simulate
              $(s’,r)$ from model, update Q
    \end{enumerate}

    The $n$ planning steps give Dyna-Q major sample-efficiency gains.
    Planning in imagination is cheap; real experience is expensive.

    \tcbitem[title={World Models \& MPC}, raster multicolumn=3]
    \textbf{World models} (Ha \& Schmidhuber, 2018) learn a compact
    latent representation of the environment; the agent plans entirely within
    the latent model.\textsuperscript{\cite[][Ch.~8]{suttonReinforcementLearningIntroduction2018}}

    \textbf{Model Predictive Control (MPC)}: at each step, optimise actions
    over a finite horizon $H$ using the model, execute only the first action,
    then re-plan. Naturally handles non-stationary environments.

    \begin{itemize}
        \item Used in robotics, process control, AlphaZero (MCTS + learned model)
        \item Horizon $H$ trades off computation vs.\ look-ahead quality
    \end{itemize}

    \end{tcbitemize}

    \subsection{Exploration vs Exploitation}
    \begin{tcbitemize}[ skin=sectionraster ]

    \tcbitem[title={The Exploration--Exploitation Trade-off}, raster multicolumn=3]
    The agent must balance:\textsuperscript{\cite[][pp.~338--339]{hurbansGrokkingArtificialIntelligence2020}}
    \begin{itemize}
        \item \textbf{Exploitation}: select $\arg\max_a Q(s,a)$ to maximise
              known reward
        \item \textbf{Exploration}: try other actions to discover potentially
              better rewards
    \end{itemize}

    Purely greedy agents get stuck in local optima; purely random agents
    never converge. Every RL algorithm must address this tension.

    \tcbitem[title=$\varepsilon$-Greedy Strategy, raster multicolumn=3]
    The simplest trade-off
    strategy:\textsuperscript{\cite[][p.~339]{hurbansGrokkingArtificialIntelligence2020}}

    \begin{equation}
        a_t = \begin{cases}
            \text{random action} & \text{with prob.\ } \varepsilon \\
            \arg\max_a Q(s_t,a) & \text{with prob.\ } 1{-}\varepsilon
        \end{cases}
    \end{equation}

    \textbf{Decaying $\varepsilon$}: start at $\varepsilon \approx 1.0$ (mostly
    explore) and anneal to $\varepsilon \approx 0.01$ (mostly exploit) over
    training:
    \[ \varepsilon_t = \varepsilon_{\min}
       + (\varepsilon_{\max} - \varepsilon_{\min})\,e^{-t/\tau} \]

    \tcbitem[title={Upper Confidence Bound (UCB)}, raster multicolumn=3]
    UCB selects actions that are either \emph{estimated good} or
    \emph{rarely tried}:\textsuperscript{\cite[][Ch.~2]{suttonReinforcementLearningIntroduction2018}}

    \begin{equation}
        a_t = \arg\max_a \left[ Q(s,a) + c\sqrt{\frac{\ln t}{N_t(s,a)}} \right]
    \end{equation}

    $N_t(s,a)$ = times action $a$ was taken from state $s$; $c$ controls
    exploration strength. The confidence bonus shrinks as an action is tried
    more, automatically directing attention to underexplored options.
    \textbf{Optimism in the face of uncertainty.}

    \tcbitem[title=Curiosity-Driven Exploration, raster multicolumn=3]
    When extrinsic rewards are sparse, \textbf{intrinsic motivation} sustains
    exploration. A \emph{forward model} $f(s_t, a_t) \to \hat{s}_{t+1}$ is
    learned; its prediction error provides an intrinsic
    reward:\textsuperscript{\cite[][Ch.~2]{suttonReinforcementLearningIntroduction2018}}

    \begin{equation}
        r^{\mathrm{int}}_t = \bigl\lVert\hat{s}_{t+1} - s_{t+1}\bigr\rVert^2
    \end{equation}

    High error signals novel, informative states the model has not yet learned.
    The agent is intrinsically \emph{curious} about regions where its world
    model is poor. Successfully applied in Montezuma’s Revenge and robotic
    manipulation.

    \end{tcbitemize}


    % -------------------------------------------------------
    \section{Inference and Causality}

    \subsection{Statistical Inference}
    \begin{tcbitemize}[ skin=sectionraster ]

    \tcbitem[title=Bayesian Inference, raster multicolumn=3]
    Bayesian inference treats model parameters $\theta$ as random variables
    and updates beliefs upon observing data
    $\mathcal{D}$:\textsuperscript{\cite[][pp.~14--19]{barberBayesianReasoningMachine2012}}

    \begin{equation}
        p(\theta|\mathcal{D}) = \frac{p(\mathcal{D}|\theta)\,p(\theta)}{p(\mathcal{D})}
    \end{equation}

    \begin{itemize}
        \item \textbf{Prior} $p(\theta)$: belief before seeing data
        \item \textbf{Likelihood} $p(\mathcal{D}|\theta)$: how probable is the data?
        \item \textbf{Posterior} $p(\theta|\mathcal{D})$: updated belief
        \item \textbf{Evidence} $p(\mathcal{D}) = \int p(\mathcal{D}|\theta)\,p(\theta)\,d\theta$
    \end{itemize}

    Posterior $\propto$ Likelihood $\times$ Prior.
    The MAP estimate is $\hat{\theta} = \arg\max_\theta p(\theta|\mathcal{D})$.

    \tcbitem[title={Bayesian vs Frequentist}, raster multicolumn=3]
    \begin{itemize}
        \item \textbf{Frequentist}: $\theta$ is a fixed unknown constant;
              probability refers to long-run frequencies; provides point
              estimates and confidence intervals
        \item \textbf{Bayesian}: $\theta$ is a random variable with a
              distribution; probability reflects degree of belief; provides
              full posterior distribution
    \end{itemize}

    Bayesian inference is naturally sequential: today’s posterior becomes
    tomorrow’s prior as new data
    arrives.\textsuperscript{\cite[][Ch.~1]{downeyThinkBayes2016}}
    Ideal for online learning and decision-making under uncertainty.

    \tcbitem[title=Bayesian Networks, raster multicolumn=6]
    A \textbf{Bayesian Network} (belief network) is a DAG where each node is
    a random variable $x_i$ with a conditional probability table (CPT)
    $p(x_i | \mathrm{pa}(x_i))$. The joint distribution
    factorises:\textsuperscript{\cite[][pp.~33--40]{barberBayesianReasoningMachine2012}}

    \begin{equation}
        p(x_1, \ldots, x_D) = \prod_{i=1}^{D} p\bigl(x_i \mid \mathrm{pa}(x_i)\bigr)
    \end{equation}

    \textbf{Wet Grass example}: $R$ (rain), $S$ (sprinkler), $T$ (Tracey’s
    grass), $J$ (Jack’s grass). With $p(T|J,R,S)=p(T|R,S)$ and $p(J|R,S)=p(J|R)$:
    \[p(T,J,R,S) = p(T|R,S)\,p(J|R)\,p(R)\,p(S)\]
    This compresses 15 parameters to 8. The graph encodes sets of conditional
    independence relations through d-separation and the local Markov property;
    an absent edge should not be read as one unconditional independence claim.
    \tcblower
    \textbf{Explaining away}: in $A \to C \leftarrow B$, causes $A$ and $B$
    are independent a priori. Conditioning on the effect $C$ makes them
    dependent: observing one cause ``explains away’’ the other.
    This is the collider effect.

    \tcbitem[title=Probabilistic Modelling, raster multicolumn=3]
    A probabilistic model defines a joint distribution over observed variables
    $\mathbf{x}$ and latent variables
    $\mathbf{z}$:\textsuperscript{\cite[][Ch.~1]{barberBayesianReasoningMachine2012}}

    \begin{equation}
        p(\mathbf{x}, \mathbf{z} | \theta) =
            p(\mathbf{x} | \mathbf{z}, \theta)\,p(\mathbf{z}|\theta)
    \end{equation}

    \textbf{Inference}: compute $p(\mathbf{z}|\mathbf{x},\theta)$ given observations.\\
    \textbf{Learning}: find $\theta$ maximising the marginal likelihood
    $p(\mathbf{x}|\theta) = \int p(\mathbf{x}|\mathbf{z},\theta)\,p(\mathbf{z}|\theta)\,d\mathbf{z}$.

    The integral is often intractable, motivating approximate methods:
    MCMC and variational inference.

    \tcbitem[title=Conditional Independence, raster multicolumn=3]
    $X$ and $Y$ are \textbf{conditionally independent} given $Z$, written
    $X \perp\!\!\!\perp Y \mid Z$,
    if:\textsuperscript{\cite[][p.~8]{barberBayesianReasoningMachine2012}}

    \begin{equation}
        p(X, Y | Z) = p(X|Z)\,p(Y|Z)
    \end{equation}

    Equivalently, $p(X|Y,Z) = p(X|Z)$: once $Z$ is known, $Y$ contains no
    additional information about $X$.

    Every absent edge in a Bayesian Network encodes exactly one conditional
    independence assumption, compressing the representation of the joint
    distribution.

    \end{tcbitemize}

    \subsection{Introduction to Causality}
    \begin{tcbitemize}[ skin=sectionraster ]

    \tcbitem[title={Correlation vs Causation}, raster multicolumn=3]
    \textbf{Correlation} is a symmetric statistical association;
    $\mathrm{Cov}(X,Y) \neq 0$ says nothing about direction or mechanism.

    \textbf{Causation} implies a directional mechanism: intervening on $X$
    forces a change in $Y$.

    Classic example: ice cream sales and drowning rates are correlated ---
    both are driven by a common cause (summer heat). Controlling for
    temperature removes the association. Formally: $p(Y|X) \neq p(Y|\,\mathrm{do}(X))$
    whenever confounders
    exist.\textsuperscript{\cite[][Ch.~1]{pearlBookWhyNew2018}}

    \tcbitem[title=Granger Causality, raster multicolumn=3]
    \textbf{Granger causality} (Granger, 1969): $X$ Granger-causes $Y$ if past
    values of $X$ provide statistically significant predictive information
    about $Y$, beyond what past $Y$ alone
    provides:\textsuperscript{\cite[][Ch.~2]{nessCausalAI2025}}

    \[ Y_t = f(Y_{t-1}, \ldots, X_{t-1}, \ldots) + \varepsilon_t \]

    Test: does adding $X_{t-k}$ significantly reduce prediction error for $Y_t$?

    \textbf{Limitation}: captures temporal precedence, not structural
    causation. Two variables driven by a hidden common cause can each appear
    to Granger-cause the other.

    \tcbitem[title={Directed Acyclic Graphs (DAGs)}, raster multicolumn=3]
    A \textbf{DAG} encodes causal structure: directed edges $X \to Y$ mean
    ``$X$ directly causes $Y$’’. Acyclicity prevents circular
    causation.\textsuperscript{\cite[][pp.~25--27]{barberBayesianReasoningMachine2012}}

    DAG terminology:
    \begin{itemize}
        \item \textbf{Parent}: direct cause, $\mathrm{pa}(Y)$
        \item \textbf{Ancestor}: any cause along a directed path
        \item \textbf{Descendant}: any variable causally downstream
        \item \textbf{Markov blanket}: parents + children + parents of children;
              screens $X$ from all other variables
    \end{itemize}

    In causal DAGs, ancestral order corresponds to temporal order.

    \tcbitem[title={Fork, Chain, and Collider}, raster multicolumn=3]
    Three elementary structures determine how information
    flows:\textsuperscript{\cite[][pp.~43--45]{barberBayesianReasoningMachine2012}}

    \begin{itemize}
        \item \textbf{Fork} (common cause): $A \leftarrow C \rightarrow B$.
              $C$ is a confounder; $A \perp\!\!\!\perp B \mid C$ but not marginally.
        \item \textbf{Chain} (mediation): $A \rightarrow C \rightarrow B$.
              $A \perp\!\!\!\perp B \mid C$ but not marginally.
        \item \textbf{Collider}: $A \rightarrow C \leftarrow B$.
              $A \perp\!\!\!\perp B$ marginally, but $A \not\!\perp\!\!\!\perp B \mid C$.
              Conditioning opens the path.
    \end{itemize}

    Conditioning on a collider (or its descendant) creates a \emph{spurious}
    association between its parents.

    \tcbitem[title=D-Separation, raster multicolumn=6]
    \textbf{D-separation} (Pearl, 1988) is the formal criterion for reading
    conditional independence from a
    DAG.\textsuperscript{\cite[][pp.~45--46]{barberBayesianReasoningMachine2012}}

    A path $U$ between $X$ and $Y$ is \textbf{blocked} by set $\mathcal{Z}$ if:
    \begin{enumerate}
        \item There is a \emph{non-collider} $w$ on $U$ with $w \in \mathcal{Z}$, OR
        \item There is a \emph{collider} $w$ on $U$ such that neither $w$ nor
              any descendant of $w$ is in $\mathcal{Z}$.
    \end{enumerate}

    $X$ and $Y$ are \textbf{d-separated} by $\mathcal{Z}$ if every path between
    them is blocked. D-separation implies $X \perp\!\!\!\perp Y \mid \mathcal{Z}$
    in any distribution consistent with the DAG structure.
    \tcblower
    \textbf{Intuition}: d-separation tells us which conditioning sets render
    variables independent. It is the key tool for identifying valid adjustment
    sets without running experiments.

    \end{tcbitemize}

    \subsection{Interventions}
    \begin{tcbitemize}[ skin=sectionraster ]

    \tcbitem[title={Seeing vs Doing: The Do-Operator}, raster multicolumn=3]
    Pearl’s \textbf{do-operator} distinguishes observation from
    intervention:\textsuperscript{\cite[][pp.~50--51]{barberBayesianReasoningMachine2012}}

    \begin{itemize}
        \item \textbf{Seeing}: $p(Y \mid X{=}x)$ --- conditioning on
              \emph{observing} $X{=}x$; confounders still act
        \item \textbf{Doing}: $p(Y \mid \mathrm{do}(X{=}x))$ --- \emph{forcing}
              $X{=}x$ by intervention; removes all incoming arrows to $X$
    \end{itemize}

    Intervening creates modified graph $\mathcal{G}_{\overline{X}}$ in which
    all edges \emph{into} $X$ are cut. Post-intervention: only $p(X_j | \mathrm{pa}(X_j))$
    for non-intervened variables remains; the intervened variable’s CPT is replaced
    by the unit distribution at the set value.

    \tcbitem[title=Confounders, raster multicolumn=3]
    A \textbf{confounder} is a variable $C$ that causally affects both
    treatment $X$ and outcome
    $Y$:\textsuperscript{\cite[][Ch.~7]{hernanCausalInferenceWhat2020}}

    \[ C \rightarrow X,\quad C \rightarrow Y \]

    Confounders create a \emph{backdoor path} from $X$ to $Y$ that does not
    represent the causal effect of $X$. Failing to adjust for confounders
    yields biased estimates.

    \textbf{Examples}: socioeconomic status confounds education and health
    outcomes; age confounds treatment choice and recovery in clinical data.
    In a DAG, confounders are identified by backdoor paths.

    \tcbitem[title=Counterfactuals, raster multicolumn=3]
    A \textbf{counterfactual} asks: ``What would $Y$ have been if $X$ had been
    $x$, given that we actually observed $X{=}x’$?’’\textsuperscript{\cite[][Ch.~7]{pearlCausalityModelsReasoning2009}}

    Notation: $Y_{X=x}(u)$ = value of $Y$ for unit $u$ under
    $\mathrm{do}(X{=}x)$.

    Requires a \textbf{Structural Causal Model} (SCM):
    $X_i = f_i(\mathrm{pa}(X_i), U_i)$.
    Three steps:
    \begin{enumerate}
        \item \textbf{Abduction}: infer exogenous noise $U$ from evidence
        \item \textbf{Action}: modify SCM to set $X{=}x$
        \item \textbf{Prediction}: compute $Y$ in the modified model
    \end{enumerate}

    \tcbitem[title={Causal Inference vs RCTs}, raster multicolumn=3]
    A \textbf{Randomized Controlled Trial (RCT)} randomly assigns treatment,
    breaking the confounder--treatment
    link:\textsuperscript{\cite[][Ch.~1]{hernanCausalInferenceWhat2020}}

    \[ \text{Random assignment} \;\Leftrightarrow\; \mathrm{do}(X)
       \;\Rightarrow\; p(Y|\mathrm{do}(X)) = p(Y|X) \]

    RCTs are the gold standard but often costly, unethical, or impossible.

    \textbf{Observational causal inference} estimates $p(Y|\mathrm{do}(X))$
    from non-experimental data via:
    \begin{itemize}
        \item Backdoor/front-door adjustment (Pearl’s criteria)
        \item Instrumental variables, regression discontinuity,
              difference-in-differences
    \end{itemize}

    \end{tcbitemize}

    \subsection{Do-Calculus}
    \begin{tcbitemize}[ skin=sectionraster ]

    \tcbitem[title=Backdoor Criterion, raster multicolumn=3]
    A set $\mathcal{Z}$ satisfies the \textbf{backdoor criterion} for
    $(X,Y)$ if:\textsuperscript{\cite[][pp.~79--81]{pearlCausalityModelsReasoning2009}}
    \begin{enumerate}
        \item No node in $\mathcal{Z}$ is a descendant of $X$
        \item $\mathcal{Z}$ blocks every backdoor path from $X$ to $Y$
    \end{enumerate}

    If satisfied, the \textbf{adjustment formula} gives:
    \begin{equation}
        p(Y \mid \mathrm{do}(X)) = \sum_{z} p(Y \mid X, \mathcal{Z}{=}z)\,p(\mathcal{Z}{=}z)
    \end{equation}

    This allows computing causal effects from observational data by conditioning
    on the confounder set $\mathcal{Z}$.

    \tcbitem[title=Front-Door Criterion, raster multicolumn=3]
    Used when all confounders are \emph{unobserved} but a mediator $M$ is
    observed:\textsuperscript{\cite[][pp.~82--84]{pearlCausalityModelsReasoning2009}}

    Conditions: (1) $M$ intercepts all directed paths $X \to Y$; (2) no
    unblocked backdoor path from $X$ to $M$; (3) all backdoor paths from $M$
    to $Y$ are blocked by $X$.

    \begin{equation}
        p(Y|\mathrm{do}(X)) = \sum_m p(M{=}m|X)
            \sum_{x’} p(Y|M{=}m, X{=}x’)\,p(X{=}x’)
    \end{equation}

    \tcbitem[title={Three Rules of Do-Calculus}, raster multicolumn=6]
    Pearl’s \textbf{do-calculus} provides three rules that together are
    \emph{complete}: any identifiable causal query can be derived from
    observational data + DAG using these
    transformations.\textsuperscript{\cite[][pp.~85--87]{pearlCausalityModelsReasoning2009}}

    Let $\mathcal{G}_{\overline{X}}$ denote graph with arrows into $X$ removed;
    $\mathcal{G}_{\underline{X}}$ graph with arrows out of $X$ removed.

    \begin{itemize}
        \item \textbf{Rule 1 (Observation)} if $(Y \perp\!\!\!\perp Z \mid X,W)_{\mathcal{G}_{\overline{X}}}$:
            $p(y|\mathrm{do}(x),z,w) = p(y|\mathrm{do}(x),w)$
        \item \textbf{Rule 2 (Action/observation exchange)} if
            $(Y \perp\!\!\!\perp Z \mid X,W)_{\mathcal{G}_{\overline{X}\underline{Z}}}$:
            $p(y|\mathrm{do}(x),\mathrm{do}(z),w) = p(y|\mathrm{do}(x),z,w)$
        \item \textbf{Rule 3 (Action deletion)} if
            $(Y \perp\!\!\!\perp Z \mid X,W)_{\mathcal{G}_{\overline{X},\overline{Z(W)}}}$:
            $p(y|\mathrm{do}(x),\mathrm{do}(z),w) = p(y|\mathrm{do}(x),w)$
    \end{itemize}

    \end{tcbitemize}

    \subsection{Fallacies}
    \begin{tcbitemize}[ skin=sectionraster ]

    \tcbitem[title=Simpson’s Paradox, raster multicolumn=3]
    A statistical trend in \emph{combined} data reverses in every
    subgroup.\textsuperscript{\cite[][pp.~49--50]{barberBayesianReasoningMachine2012}}

    \textbf{Medical trial example}: A drug appears beneficial overall
    (recovery: 50\% drug vs.\ 40\% no-drug) but is \emph{harmful} in each
    subgroup (males: 60\% vs.\ 70\%; females: 20\% vs.\ 30\%).

    \textbf{Resolution}: The paradox conflates $p(R|D)$ (observational)
    with $p(R|\mathrm{do}(D))$ (interventional). Gender is a confounder
    affecting both drug choice and recovery. Applying the backdoor
    adjustment reveals the drug has no net benefit.

    \tcbitem[title={Collider Bias \& Berkson’s Paradox}, raster multicolumn=3]
    \textbf{Collider bias}: conditioning on a common effect of two variables
    creates a spurious association between
    them.\textsuperscript{\cite[][p.~43]{barberBayesianReasoningMachine2012}}

    Structure: $A \to C \leftarrow B$. Marginally $A \perp\!\!\!\perp B$, but
    $A \not\!\perp\!\!\!\perp B \mid C$.

    \textbf{Berkson’s Paradox} (1946): In hospital studies, diseases $A$ and
    $B$ appear negatively correlated because hospitalisation is a collider.
    Patients with $A$ seem less likely to also have $B$ --- but only because
    both independently cause hospitalisation. In the general population, $A$
    and $B$ are independent.

    \tcbitem[title=Mediation Fallacy, raster multicolumn=3]
    The \textbf{mediation fallacy} occurs when a researcher adjusts for a
    mediator $M$ on the causal path $X \to M \to Y$ while intending to
    estimate the total effect of $X$ on
    $Y$.\textsuperscript{\cite[][Ch.~9]{pearlCausalityModelsReasoning2009}}

    Adjusting for $M$:
    \begin{itemize}
        \item \emph{Blocks} the indirect pathway $X \to M \to Y$
        \item The estimate captures only the \emph{direct} effect $X \to Y$,
              not the total effect
        \item If the goal is the total effect, $M$ must \emph{not} be
              conditioned on
    \end{itemize}

    Correct mediation analysis separates direct and indirect effects using
    counterfactual definitions.

    \tcbitem[title={Missing Values: Causal View}, raster multicolumn=3]
    Standard imputation (mean/median/model) assumes data is Missing At
    Random (MAR). The causal view requires modelling the
    \textbf{missingness mechanism}:\textsuperscript{\cite[][Ch.~8]{nessCausalAI2025}}

    \begin{itemize}
        \item \textbf{MCAR} (Missing Completely At Random): missingness
              independent of all variables --- safe to impute
        \item \textbf{MAR} (Missing At Random): missingness depends on
              observed variables --- adjust for them
        \item \textbf{MNAR} (Missing Not At Random): missingness depends
              on the missing value itself --- selection bias; model the
              mechanism explicitly
    \end{itemize}

    If the variable with missing values is a collider descendant, naive
    imputation can introduce collider bias. The DAG determines the
    correct imputation strategy.

    \end{tcbitemize}

    % -------------------------------------------------------
    \section{Further Reading}
    \begin{tcbitemize}[ skin=sectionraster ]

    \tcbitem[title=Recommended Sources for Reinforcement Learning and Causality, raster multicolumn=6]
    \begin{itemize}
        \item \fullcite{suttonReinforcementLearningIntroduction2018} --- The definitive RL textbook: MDPs, TD learning, policy gradients, function approximation, and applications.
        \item \fullcite{hurbansGrokkingArtificialIntelligence2020} --- Accessible introduction to RL with visual examples, Q-learning, and exploration strategies.
        \item \fullcite{barberBayesianReasoningMachine2012} --- Bayesian inference, graphical models, MDPs, and value iteration with rigorous probabilistic treatment.
        \item \fullcite{downeyThinkBayes2016} --- Practical Bayesian statistics with Python: conjugate priors, MCMC, and sequential updating.
        \item \fullcite{pearlCausalityModelsReasoning2009} --- The foundational text on structural causal models, do-calculus, and causal identification.
        \item \fullcite{pearlBookWhyNew2018} --- Popular-science introduction to the ``ladder of causation'' and the do-operator.
        \item \fullcite{pearlCausalInferenceStatistics2016} --- Concise primer on causal inference: DAGs, adjustment, frontdoor/backdoor criteria.
        \item \fullcite{hernanCausalInferenceWhat2020} --- Practical causal inference for health sciences: RCTs, observational studies, time-varying treatments.
        \item \fullcite{nessCausalAI2025} --- Modern perspective on causal AI, including Granger causality, causal discovery, and missing data.
    \end{itemize}

    \end{tcbitemize}
